% Part: first-order-logic
% Chapter: tableaux
% Section: proof-theoretic-notions

\documentclass[../../../include/open-logic-section]{subfiles}

\begin{document}

\iftag{FOL}
      {\olfileid{fol}{tab}{ptn}}
      {\olfileid{pl}{tab}{ptn}}

\olsection{প্রমাণতাত্ত্বিক ধারণা}

\begin{editorial}
এই অংশে ট্যাবলোর প্রমাণযোগ্যতা-সম্পর্ক ও সঙ্গতির সংজ্ঞাগুলি একত্র
করা হয়েছে।
\end{editorial}

\begin{explain}
বৈধতা, যৌক্তিক ফলশ্রুতি ও পরিতৃপ্তিযোগ্যতার মতো কয়েকটি গুরুত্বপূর্ণ
অর্থতাত্ত্বিক ধারণা সংজ্ঞায়িত করেছি। এখন তাদের অনুরূপ
\emph{প্রমাণতাত্ত্বিক ধারণা} সংজ্ঞায়িত করি। এগুলি !!{structure}s-এ
!!{sentence}s-এর পরিতৃপ্তির ভিত্তিতে নয়, কয়েকটি বিশেষ বদ্ধ
!!{tableau}x-এর অস্তিত্বের ভিত্তিতে সংজ্ঞায়িত। এই দুই ধরনের ধারণা যে
মিলে যায়, তা একটি গুরুত্বপূর্ণ আবিষ্কার। \emph{বিশুদ্ধতা} ও
\emph{পূর্ণতা উপপাদ্য} এই মিলই প্রতিষ্ঠা করে।
\end{explain}


\begin{defn}[উপপাদ্য]
কোনো !!{sentence}~$!A$-এর জন্য $\sFmla{\False}{!A}$-এর একটি বদ্ধ
!!{tableau} থাকলে $!A$-কে \emph{উপপাদ্য} বলা হয়। সেটি উপপাদ্য হলে
$\Proves !A$ এবং উপপাদ্য না হলে $\Proves/ !A$ লিখি।
\end{defn}

\begin{defn}[!!^{derivability}]
!!^a{sentence}~$!A$ কোনো !!{sentence}s-এর সমষ্টি~$\Gamma$ থেকে
\emph{!!{derivable}} হলে, অর্থাৎ $\Gamma \Proves !A$ হবে তখনই, যখন
কোনো সসীম সমষ্টি $\{!B_1, \dots, !B_n\} \subseteq \Gamma$ এবং
নিচের সমষ্টিটির জন্য একটি বদ্ধ !!{tableau} থাকে:
\[
\{
  \sFmla{\False}{!A},
  \sFmla{\True}{!B_1}, \dots,
  \sFmla{\True}{!B_n}
  \}.
\]
$!A$ যদি $\Gamma$ থেকে !!{derivable} না হয়, তবে
$\Gamma \Proves/ !A$ লিখি।
\end{defn}

\begin{defn}[সঙ্গতি]
!!{sentence}s-এর কোনো সমষ্টি~$\Gamma$ \emph{অসঙ্গত} হবে তখনই, যখন
কোনো সসীম সমষ্টি $\{!B_1, \dots, !B_n\} \subseteq \Gamma$ এবং
নিচের সমষ্টিটির জন্য একটি বদ্ধ !!{tableau} থাকে:
\[
\{
  \sFmla{\True}{!B_1}, \dots,
  \sFmla{\True}{!B_n}
  \}.
\]
$\Gamma$ অসঙ্গত না হলে তাকে \emph{সঙ্গত} বলি।
\end{defn}

\begin{prop}[প্রতিবিম্ব ধর্ম]
\ollabel{prop:reflexivity}
$!A \in \Gamma$ হলে $\Gamma \Proves !A$।
\end{prop}

\begin{proof}
$!A \in \Gamma$ হলে $\{!A\}$ হলো~$\Gamma$-র সসীম উপসমষ্টি এবং
!!{tableau}
\begin{oltableau}
  [\sFmla{\False}{\formula{A}}, just = \TAss
    [\sFmla{\True}{\formula{A}}, just = \TAss,close]
  ]
\end{oltableau}
বদ্ধ।
\end{proof}


\begin{prop}[একঘেয়েতা]
\ollabel{prop:monotonicity}
$\Gamma \subseteq \Delta$ এবং $\Gamma \Proves !A$ হলে
$\Delta \Proves !A$।
\end{prop}

\begin{proof}
~$\Gamma$-র যেকোনো সসীম উপসমষ্টি~$\Delta$-রও সসীম উপসমষ্টি।
\end{proof}

\begin{prop}[পরিযায়িতা]
\ollabel{prop:transitivity}
$\Gamma \Proves !A$ এবং $\{!A\} \cup
\Delta \Proves !B$ হলে $\Gamma \cup \Delta \Proves !B$।
\end{prop}

\begin{proof}
$\{!A\} \cup \Delta \Proves !B$ হলে
$\Delta_0=\{!C_1,\dots,!C_n\}\subseteq\Delta$ এমন একটি সসীম
উপসমষ্টি আছে, যার জন্য
\begin{align*}
\{\sFmla{\False}{!B}, & \sFmla{\True}{!A}, \sFmla{\True}{!C_1},
\dots, \sFmla{\True}{!C_n}\}
\intertext{-এর একটি বদ্ধ !!{tableau} থাকে। আর $\Gamma \Proves !A$
হলে $\{!D_1,\dots,!D_m\}\subseteq\Gamma$ এমন একটি সসীম উপসমষ্টি
আছে, যার জন্য}
\{\sFmla{\False}{!A}, & \sFmla{\True}{!D_1},
\dots, \sFmla{\True}{!D_m}\}
\end{align*}
-এর একটি বদ্ধ !!{tableau} থাকে।
% BN-SRC-049 TARGET-CORRECTION-BEGIN
\par\smallskip\noindent\textnormal{\small\textbf{উৎস-সংশোধন (BN-SRC-049):} হিমায়িত ইংরেজি গদ্যে পৃথক বাক্যগুলির তালিকা $!D_1,\dots,!D_m$-কেই $\Gamma$-র উপসমষ্টি বলা হয়েছে। আগের $\Delta_0$-সংজ্ঞা ও পরের ট্যাবলো-অনুমিতির সঙ্গে মিলিয়ে বাংলা লক্ষ্যে সমষ্টি-বন্ধনী পুনরুদ্ধার করে $\{!D_1,\dots,!D_m\}\subseteq\Gamma$ লেখা হয়েছে।} % BN-SRC-049 TARGET-CORRECTION-END

এবার নিচের অনুমিতিযুক্ত !!{tableau} বিবেচনা করো:
\[
\sFmla{\False}{!B},
\sFmla{\True}{!C_1}, \dots, \sFmla{\True}{!C_n},
\sFmla{\True}{!D_1}, \dots, \sFmla{\True}{!D_m}.
\]
$!A$-তে \Cut{} বিধি প্রয়োগ করো। এতে দুটি শাখা তৈরি হয়—একটিতে
$\sFmla{\True}{!A}$ এবং অন্যটিতে $\sFmla{\False}{!A}$। অতএব প্রথম
শাখায় নিচের সবগুলি পাওয়া যায়:
\[
\{\sFmla{\False}{!B}, \sFmla{\True}{!A},
\sFmla{\True}{!C_1}, \dots, \sFmla{\True}{!C_n}\}
\]
এই অনুমিতিগুলির জন্য বদ্ধ !!{tableau} আছে বলে সেটি সেই শাখায় জুড়ে
দেওয়া যায়; $\sFmla{\True}{!A}$ দিয়ে যাওয়া প্রতিটি শাখা বদ্ধ হয়।
অন্য শাখায় নিচের সবগুলি পাওয়া যায়:
\[
\{\sFmla{\False}{!A}, \sFmla{\True}{!D_1}, \dots,
\sFmla{\True}{!D_m}\}
\]
তাই অন্য পাশটিও সম্পূর্ণ করে একটি বদ্ধ !!{tableau} পাওয়া যায়।
এতে $\Gamma \cup \Delta \Proves !B$ প্রমাণিত হলো।
\end{proof}

লক্ষ করো, এর বিশেষ ক্ষেত্র হিসেবে $\Gamma \Proves !A$ এবং
$!A \Proves !B$ হলে $\Gamma \Proves !B$। আবার,
$!A_1,\dots,!A_n \Proves !B$ এবং প্রতিটি~$i$-এর জন্য
$\Gamma \Proves !A_i$ হলে $\Gamma \Proves !B$।

\begin{prop}
\ollabel{prop:incons}
$\Gamma$ অসঙ্গত হবে তখনই এবং কেবল তখনই, যখন প্রতিটি
!!{sentence}~$!A$-এর জন্য $\Gamma \Proves !A$।
\end{prop}

\begin{proof}
অনুশীলন।
\end{proof}

\tagprob{FOL}
\begin{prob}
\olref[fol][tab][ptn]{prop:incons} প্রমাণ করো।
\end{prob}
\tagendprob

\tagprob{notFOL}
\begin{prob}
\olref[pl][tab][ptn]{prop:incons} প্রমাণ করো।
\end{prob}
\tagendprob

\begin{prop}[সংহতি]
  \ollabel{prop:proves-compact}
  \begin{enumerate}
  \item $\Gamma \Proves !A$ হলে $\Gamma_0\subseteq\Gamma$ এমন একটি
    সসীম উপসমষ্টি আছে যেন $\Gamma_0 \Proves !A$।
  \item ~$\Gamma$-র প্রতিটি সসীম উপসমষ্টি সঙ্গত হলে $\Gamma$ সঙ্গত।
  \end{enumerate}
\end{prop}

\begin{proof}
  \begin{enumerate}
    \item $\Gamma \Proves !A$ হলে
      $\Gamma_0=\{!B_1,\dots,!B_n\}$ এমন একটি সসীম উপসমষ্টি এবং
      নিচের সমষ্টিটির জন্য একটি বদ্ধ !!{tableau} আছে:
      \[
      \{\sFmla{\False}{!A}, \sFmla{\True}{!B_1}, \dots, \sFmla{\True}{!B_n}\}
      \]
      এই !!{tableau}-টিই $\Gamma_0 \Proves !A$ দেখায়।
    \item $\Gamma$ অসঙ্গত হলে কোনো সসীম উপসমষ্টি
      $\Gamma_0=\{!B_1,\dots,!B_n\}$-এর জন্য নিচের সমষ্টিটির একটি
      বদ্ধ !!{tableau} আছে:
      \[
      \{\sFmla{\True}{!B_1}, \dots, \sFmla{\True}{!B_n}\}
      \]
      এই বদ্ধ !!{tableau}-টি দেখায় যে $\Gamma_0$ অসঙ্গত।
  \end{enumerate}
\end{proof}

\end{document}
